% =====================================================================
%  Chapter 1 — also serves as a showcase of every book element.
%  Replace the sample content once the look is settled.
% =====================================================================
\chapter{Chapter 1 Title}
\label{ch:01}

This opening chapter demonstrates each element the book supports.
Body text is set in EB Garamond with old-style figures such as 1984 and
2026.\footnote{This is an ordinary footnote. It appears at the bottom of
the page and is numbered per chapter.} A second footnote shows how
numbering continues.\footnote{Footnotes can run to several lines without
trouble; the hanging indent keeps the number clear of the text so that
the note remains easy to scan.}

\aside{This is a margin aside. Short notes sit in the wide outer margin,
beside the line where they are called.}
Margin asides are created with the \verb|\aside| command. They always
appear in the outer margin, on the right of odd pages and the left of
even pages.

\section{Charts}

Charts are drawn with \textsf{pgfplots}, reading their numbers from CSV
files in the \texttt{data/} folder, as in Figure~\ref{fig:sample-chart}.

\begin{figure}[htbp]
  \centering
  % Sample pgfplots chart reading data/sample.csv
\begin{tikzpicture}
  \begin{axis}[
    width=\linewidth, height=0.65\linewidth,
    xlabel={Year}, ylabel={Value},
    xticklabel style={/pgf/number format/1000 sep=}, xtick distance=4,
    grid=major, legend pos=north west, legend style={draw=none, font=\footnotesize},
  ]
    \addplot table[x=year, y=seriesA, col sep=comma] {data/sample.csv};
    \addplot table[x=year, y=seriesB, col sep=comma] {data/sample.csv};
    \legend{Series A, Series B}
  \end{axis}
\end{tikzpicture}

  \caption{A sample chart drawn from \texttt{data/sample.csv}.}
  \label{fig:sample-chart}
\end{figure}

\section{Diagrams}

Diagrams are drawn with Ti\emph{k}Z. Figure~\ref{fig:sample-diagram}
uses the book's colour palette.

\begin{figure}[htbp]
  \centering
  % Sample TikZ diagram
\begin{tikzpicture}[
    node distance=6mm and 8mm,
    box/.style={draw=taoaccent, thick, rounded corners=2pt, fill=taopaper,
                minimum width=18mm, minimum height=8mm, font=\small},
    arrow/.style={-{Stealth[length=2mm]}, taoink!70, thick},
  ]
  \node[box] (sand) {Sand};
  \node[box, right=of sand] (silicon) {Silicon};
  \node[box, right=of silicon] (chip) {Chip};
  \node[box, below=of silicon] (thought) {Thought};
  \draw[arrow] (sand) -- (silicon);
  \draw[arrow] (silicon) -- (chip);
  \draw[arrow] (chip) |- (thought);
  \draw[arrow, dashed] (thought) -| (sand);
\end{tikzpicture}

  \caption{A sample Ti\emph{k}Z diagram.}
  \label{fig:sample-diagram}
\end{figure}

\section{Sidebars}

A boxed sidebar sets material apart from the main argument.

\begin{sidebox}[A Sidebar Title]
  Sidebars break across pages when necessary and accept an optional
  title. This note uses \verb|\boxfootnote|, so its note goes to
  the bottom of the page.\boxfootnote{A footnote called from inside
  a sidebox, numbered with the chapter's other footnotes.}
\end{sidebox}

\begin{sidebox}
  A sidebar may also be untitled.
\end{sidebox}

\section{Wide material}

When a chart needs more room, the \texttt{wide} environment lets it
extend into the outer margin.

\begin{figure}[htbp]
  \begin{wide}
    \centering
    % Sample pgfplots chart reading data/sample.csv
\begin{tikzpicture}
  \begin{axis}[
    width=\linewidth, height=0.65\linewidth,
    xlabel={Year}, ylabel={Value},
    xticklabel style={/pgf/number format/1000 sep=}, xtick distance=4,
    grid=major, legend pos=north west, legend style={draw=none, font=\footnotesize},
  ]
    \addplot table[x=year, y=seriesA, col sep=comma] {data/sample.csv};
    \addplot table[x=year, y=seriesB, col sep=comma] {data/sample.csv};
    \legend{Series A, Series B}
  \end{axis}
\end{tikzpicture}

    \caption{The same chart set full width.}
    \label{fig:sample-chart-wide}
  \end{wide}
\end{figure}

After the figures, footnote numbering carries on as
expected.\footnote{The fourth footnote of this chapter.}
